% Môn: Vật lý
% Lớp: 12
% Chương: Chương I. Vật lí nhiệt
% Bài: L12C1 Bài 1. Cấu trúc của chất. Sự chuyển thể
% Loại: Lý thuyết
% Bộ sách: Kết nối tri thức với cuộc sống
% File riêng, không trộn vào de.tex
% Định nghĩa lệnh Lý thuyết — dùng khi biên dịch lt.tex bằng TeX.
% App web đọc lt.tex trực tiếp (bỏ \input file này).
% Trong lt.tex, từ thư mục bài:
%   % Định nghĩa lệnh Lý thuyết — dùng khi biên dịch lt.tex bằng TeX.
% App web đọc lt.tex trực tiếp (bỏ \input file này).
% Trong lt.tex, từ thư mục bài:
%   % Định nghĩa lệnh Lý thuyết — dùng khi biên dịch lt.tex bằng TeX.
% App web đọc lt.tex trực tiếp (bỏ \input file này).
% Trong lt.tex, từ thư mục bài:
%   \input{../../../../_lenh/lythuyet.tex}
% từ gốc repo:
%   \input{ngan-hang/_lenh/lythuyet.tex}
\makeatletter
\providecommand{\indam}[1]{\textbf{#1}}
\providecommand{\chuy}[1]{\par\smallskip\noindent\textit{#1}\par}
\providecommand{\luuy}[1]{\par\smallskip\noindent\textbf{Lưu ý.} #1\par}
\providecommand{\ghichu}[1]{\par\smallskip\noindent\textbf{Ghi chú.} #1\par}
\providecommand{\vidu}[1]{\par\smallskip\noindent\textbf{Ví dụ.} #1\par}
\providecommand{\Blythuyet}{\subsection*{Lý thuyết}}
% Hai cột: kiến thức | hình
\providecommand{\haicotchay}[2]{%
  \par\medskip\noindent
  \begin{minipage}[t]{0.52\linewidth}#1\end{minipage}\hfill
  \begin{minipage}[t]{0.46\linewidth}#2\end{minipage}%
  \par\medskip
}
\providecommand{\kienthuccot}[2]{\haicotchay{#1}{#2}}
% Dự phòng nếu chưa nạp graphicx / caption (không ghi đè bản thật)
\providecommand{\resizebox}[3]{#3}
\providecommand{\captionof}[2]{\par\smallskip\noindent\textit{#2}\par}
\newenvironment{hoatdong}[1]{%
  \par\medskip\noindent\textbf{Hoạt động.} #1\par\smallskip
}{\par\medskip}
\newenvironment{emcobiet}{%
  \par\medskip\noindent\textbf{Em có biết}\par\smallskip
}{\par\medskip}
\newenvironment{emdahoc}{%
  \par\medskip\noindent\textbf{Em đã học}\par\smallskip
}{\par\medskip}
\newenvironment{emcothe}{%
  \par\medskip\noindent\textbf{Em có thể}\par\smallskip
}{\par\medskip}
\newenvironment{kienthuc}{%
  \par\medskip\noindent\textbf{Kiến thức cốt lõi}\par\smallskip
}{\par\medskip}
\newenvironment{traloi}{%
  \par\smallskip\noindent\textit{Trả lời / ghi chú của em}\par
  \vspace{1.2em}%
}{\par\medskip}
\makeatother

% từ gốc repo:
%   % Định nghĩa lệnh Lý thuyết — dùng khi biên dịch lt.tex bằng TeX.
% App web đọc lt.tex trực tiếp (bỏ \input file này).
% Trong lt.tex, từ thư mục bài:
%   \input{../../../../_lenh/lythuyet.tex}
% từ gốc repo:
%   \input{ngan-hang/_lenh/lythuyet.tex}
\makeatletter
\providecommand{\indam}[1]{\textbf{#1}}
\providecommand{\chuy}[1]{\par\smallskip\noindent\textit{#1}\par}
\providecommand{\luuy}[1]{\par\smallskip\noindent\textbf{Lưu ý.} #1\par}
\providecommand{\ghichu}[1]{\par\smallskip\noindent\textbf{Ghi chú.} #1\par}
\providecommand{\vidu}[1]{\par\smallskip\noindent\textbf{Ví dụ.} #1\par}
\providecommand{\Blythuyet}{\subsection*{Lý thuyết}}
% Hai cột: kiến thức | hình
\providecommand{\haicotchay}[2]{%
  \par\medskip\noindent
  \begin{minipage}[t]{0.52\linewidth}#1\end{minipage}\hfill
  \begin{minipage}[t]{0.46\linewidth}#2\end{minipage}%
  \par\medskip
}
\providecommand{\kienthuccot}[2]{\haicotchay{#1}{#2}}
% Dự phòng nếu chưa nạp graphicx / caption (không ghi đè bản thật)
\providecommand{\resizebox}[3]{#3}
\providecommand{\captionof}[2]{\par\smallskip\noindent\textit{#2}\par}
\newenvironment{hoatdong}[1]{%
  \par\medskip\noindent\textbf{Hoạt động.} #1\par\smallskip
}{\par\medskip}
\newenvironment{emcobiet}{%
  \par\medskip\noindent\textbf{Em có biết}\par\smallskip
}{\par\medskip}
\newenvironment{emdahoc}{%
  \par\medskip\noindent\textbf{Em đã học}\par\smallskip
}{\par\medskip}
\newenvironment{emcothe}{%
  \par\medskip\noindent\textbf{Em có thể}\par\smallskip
}{\par\medskip}
\newenvironment{kienthuc}{%
  \par\medskip\noindent\textbf{Kiến thức cốt lõi}\par\smallskip
}{\par\medskip}
\newenvironment{traloi}{%
  \par\smallskip\noindent\textit{Trả lời / ghi chú của em}\par
  \vspace{1.2em}%
}{\par\medskip}
\makeatother

\makeatletter
\providecommand{\indam}[1]{\textbf{#1}}
\providecommand{\chuy}[1]{\par\smallskip\noindent\textit{#1}\par}
\providecommand{\luuy}[1]{\par\smallskip\noindent\textbf{Lưu ý.} #1\par}
\providecommand{\ghichu}[1]{\par\smallskip\noindent\textbf{Ghi chú.} #1\par}
\providecommand{\vidu}[1]{\par\smallskip\noindent\textbf{Ví dụ.} #1\par}
\providecommand{\Blythuyet}{\subsection*{Lý thuyết}}
% Hai cột: kiến thức | hình
\providecommand{\haicotchay}[2]{%
  \par\medskip\noindent
  \begin{minipage}[t]{0.52\linewidth}#1\end{minipage}\hfill
  \begin{minipage}[t]{0.46\linewidth}#2\end{minipage}%
  \par\medskip
}
\providecommand{\kienthuccot}[2]{\haicotchay{#1}{#2}}
% Dự phòng nếu chưa nạp graphicx / caption (không ghi đè bản thật)
\providecommand{\resizebox}[3]{#3}
\providecommand{\captionof}[2]{\par\smallskip\noindent\textit{#2}\par}
\newenvironment{hoatdong}[1]{%
  \par\medskip\noindent\textbf{Hoạt động.} #1\par\smallskip
}{\par\medskip}
\newenvironment{emcobiet}{%
  \par\medskip\noindent\textbf{Em có biết}\par\smallskip
}{\par\medskip}
\newenvironment{emdahoc}{%
  \par\medskip\noindent\textbf{Em đã học}\par\smallskip
}{\par\medskip}
\newenvironment{emcothe}{%
  \par\medskip\noindent\textbf{Em có thể}\par\smallskip
}{\par\medskip}
\newenvironment{kienthuc}{%
  \par\medskip\noindent\textbf{Kiến thức cốt lõi}\par\smallskip
}{\par\medskip}
\newenvironment{traloi}{%
  \par\smallskip\noindent\textit{Trả lời / ghi chú của em}\par
  \vspace{1.2em}%
}{\par\medskip}
\makeatother

% từ gốc repo:
%   % Định nghĩa lệnh Lý thuyết — dùng khi biên dịch lt.tex bằng TeX.
% App web đọc lt.tex trực tiếp (bỏ \input file này).
% Trong lt.tex, từ thư mục bài:
%   % Định nghĩa lệnh Lý thuyết — dùng khi biên dịch lt.tex bằng TeX.
% App web đọc lt.tex trực tiếp (bỏ \input file này).
% Trong lt.tex, từ thư mục bài:
%   \input{../../../../_lenh/lythuyet.tex}
% từ gốc repo:
%   \input{ngan-hang/_lenh/lythuyet.tex}
\makeatletter
\providecommand{\indam}[1]{\textbf{#1}}
\providecommand{\chuy}[1]{\par\smallskip\noindent\textit{#1}\par}
\providecommand{\luuy}[1]{\par\smallskip\noindent\textbf{Lưu ý.} #1\par}
\providecommand{\ghichu}[1]{\par\smallskip\noindent\textbf{Ghi chú.} #1\par}
\providecommand{\vidu}[1]{\par\smallskip\noindent\textbf{Ví dụ.} #1\par}
\providecommand{\Blythuyet}{\subsection*{Lý thuyết}}
% Hai cột: kiến thức | hình
\providecommand{\haicotchay}[2]{%
  \par\medskip\noindent
  \begin{minipage}[t]{0.52\linewidth}#1\end{minipage}\hfill
  \begin{minipage}[t]{0.46\linewidth}#2\end{minipage}%
  \par\medskip
}
\providecommand{\kienthuccot}[2]{\haicotchay{#1}{#2}}
% Dự phòng nếu chưa nạp graphicx / caption (không ghi đè bản thật)
\providecommand{\resizebox}[3]{#3}
\providecommand{\captionof}[2]{\par\smallskip\noindent\textit{#2}\par}
\newenvironment{hoatdong}[1]{%
  \par\medskip\noindent\textbf{Hoạt động.} #1\par\smallskip
}{\par\medskip}
\newenvironment{emcobiet}{%
  \par\medskip\noindent\textbf{Em có biết}\par\smallskip
}{\par\medskip}
\newenvironment{emdahoc}{%
  \par\medskip\noindent\textbf{Em đã học}\par\smallskip
}{\par\medskip}
\newenvironment{emcothe}{%
  \par\medskip\noindent\textbf{Em có thể}\par\smallskip
}{\par\medskip}
\newenvironment{kienthuc}{%
  \par\medskip\noindent\textbf{Kiến thức cốt lõi}\par\smallskip
}{\par\medskip}
\newenvironment{traloi}{%
  \par\smallskip\noindent\textit{Trả lời / ghi chú của em}\par
  \vspace{1.2em}%
}{\par\medskip}
\makeatother

% từ gốc repo:
%   % Định nghĩa lệnh Lý thuyết — dùng khi biên dịch lt.tex bằng TeX.
% App web đọc lt.tex trực tiếp (bỏ \input file này).
% Trong lt.tex, từ thư mục bài:
%   \input{../../../../_lenh/lythuyet.tex}
% từ gốc repo:
%   \input{ngan-hang/_lenh/lythuyet.tex}
\makeatletter
\providecommand{\indam}[1]{\textbf{#1}}
\providecommand{\chuy}[1]{\par\smallskip\noindent\textit{#1}\par}
\providecommand{\luuy}[1]{\par\smallskip\noindent\textbf{Lưu ý.} #1\par}
\providecommand{\ghichu}[1]{\par\smallskip\noindent\textbf{Ghi chú.} #1\par}
\providecommand{\vidu}[1]{\par\smallskip\noindent\textbf{Ví dụ.} #1\par}
\providecommand{\Blythuyet}{\subsection*{Lý thuyết}}
% Hai cột: kiến thức | hình
\providecommand{\haicotchay}[2]{%
  \par\medskip\noindent
  \begin{minipage}[t]{0.52\linewidth}#1\end{minipage}\hfill
  \begin{minipage}[t]{0.46\linewidth}#2\end{minipage}%
  \par\medskip
}
\providecommand{\kienthuccot}[2]{\haicotchay{#1}{#2}}
% Dự phòng nếu chưa nạp graphicx / caption (không ghi đè bản thật)
\providecommand{\resizebox}[3]{#3}
\providecommand{\captionof}[2]{\par\smallskip\noindent\textit{#2}\par}
\newenvironment{hoatdong}[1]{%
  \par\medskip\noindent\textbf{Hoạt động.} #1\par\smallskip
}{\par\medskip}
\newenvironment{emcobiet}{%
  \par\medskip\noindent\textbf{Em có biết}\par\smallskip
}{\par\medskip}
\newenvironment{emdahoc}{%
  \par\medskip\noindent\textbf{Em đã học}\par\smallskip
}{\par\medskip}
\newenvironment{emcothe}{%
  \par\medskip\noindent\textbf{Em có thể}\par\smallskip
}{\par\medskip}
\newenvironment{kienthuc}{%
  \par\medskip\noindent\textbf{Kiến thức cốt lõi}\par\smallskip
}{\par\medskip}
\newenvironment{traloi}{%
  \par\smallskip\noindent\textit{Trả lời / ghi chú của em}\par
  \vspace{1.2em}%
}{\par\medskip}
\makeatother

\makeatletter
\providecommand{\indam}[1]{\textbf{#1}}
\providecommand{\chuy}[1]{\par\smallskip\noindent\textit{#1}\par}
\providecommand{\luuy}[1]{\par\smallskip\noindent\textbf{Lưu ý.} #1\par}
\providecommand{\ghichu}[1]{\par\smallskip\noindent\textbf{Ghi chú.} #1\par}
\providecommand{\vidu}[1]{\par\smallskip\noindent\textbf{Ví dụ.} #1\par}
\providecommand{\Blythuyet}{\subsection*{Lý thuyết}}
% Hai cột: kiến thức | hình
\providecommand{\haicotchay}[2]{%
  \par\medskip\noindent
  \begin{minipage}[t]{0.52\linewidth}#1\end{minipage}\hfill
  \begin{minipage}[t]{0.46\linewidth}#2\end{minipage}%
  \par\medskip
}
\providecommand{\kienthuccot}[2]{\haicotchay{#1}{#2}}
% Dự phòng nếu chưa nạp graphicx / caption (không ghi đè bản thật)
\providecommand{\resizebox}[3]{#3}
\providecommand{\captionof}[2]{\par\smallskip\noindent\textit{#2}\par}
\newenvironment{hoatdong}[1]{%
  \par\medskip\noindent\textbf{Hoạt động.} #1\par\smallskip
}{\par\medskip}
\newenvironment{emcobiet}{%
  \par\medskip\noindent\textbf{Em có biết}\par\smallskip
}{\par\medskip}
\newenvironment{emdahoc}{%
  \par\medskip\noindent\textbf{Em đã học}\par\smallskip
}{\par\medskip}
\newenvironment{emcothe}{%
  \par\medskip\noindent\textbf{Em có thể}\par\smallskip
}{\par\medskip}
\newenvironment{kienthuc}{%
  \par\medskip\noindent\textbf{Kiến thức cốt lõi}\par\smallskip
}{\par\medskip}
\newenvironment{traloi}{%
  \par\smallskip\noindent\textit{Trả lời / ghi chú của em}\par
  \vspace{1.2em}%
}{\par\medskip}
\makeatother

\makeatletter
\providecommand{\indam}[1]{\textbf{#1}}
\providecommand{\chuy}[1]{\par\smallskip\noindent\textit{#1}\par}
\providecommand{\luuy}[1]{\par\smallskip\noindent\textbf{Lưu ý.} #1\par}
\providecommand{\ghichu}[1]{\par\smallskip\noindent\textbf{Ghi chú.} #1\par}
\providecommand{\vidu}[1]{\par\smallskip\noindent\textbf{Ví dụ.} #1\par}
\providecommand{\Blythuyet}{\subsection*{Lý thuyết}}
% Hai cột: kiến thức | hình
\providecommand{\haicotchay}[2]{%
  \par\medskip\noindent
  \begin{minipage}[t]{0.52\linewidth}#1\end{minipage}\hfill
  \begin{minipage}[t]{0.46\linewidth}#2\end{minipage}%
  \par\medskip
}
\providecommand{\kienthuccot}[2]{\haicotchay{#1}{#2}}
% Dự phòng nếu chưa nạp graphicx / caption (không ghi đè bản thật)
\providecommand{\resizebox}[3]{#3}
\providecommand{\captionof}[2]{\par\smallskip\noindent\textit{#2}\par}
\newenvironment{hoatdong}[1]{%
  \par\medskip\noindent\textbf{Hoạt động.} #1\par\smallskip
}{\par\medskip}
\newenvironment{emcobiet}{%
  \par\medskip\noindent\textbf{Em có biết}\par\smallskip
}{\par\medskip}
\newenvironment{emdahoc}{%
  \par\medskip\noindent\textbf{Em đã học}\par\smallskip
}{\par\medskip}
\newenvironment{emcothe}{%
  \par\medskip\noindent\textbf{Em có thể}\par\smallskip
}{\par\medskip}
\newenvironment{kienthuc}{%
  \par\medskip\noindent\textbf{Kiến thức cốt lõi}\par\smallskip
}{\par\medskip}
\newenvironment{traloi}{%
  \par\smallskip\noindent\textit{Trả lời / ghi chú của em}\par
  \vspace{1.2em}%
}{\par\medskip}
\makeatother


% (Không có btmau)
\subsection{Lý Thuyết}
%\Blythuyet
% Lớp: 12
% Bài: Bài 1. Cấu trúc của chất. Sự chuyển thể
% Bộ sách/tài liệu: Kết nối tri thức với cuộc sống
%**Danh sách các dạng bài tập – Bài 1. Sự Chuyển Thể**
%1. Mô hình động học phân tử về cấu tạo chất và chuyển động Brown
%2. Phân biệt đặc điểm cấu trúc phân tử của chất ở thể rắn, lỏng, khí
%3. Phân biệt chất rắn kết tinh và chất rắn vô định hình
%4. Giải thích cơ chế chuyển thể dựa trên mô hình động học phân tử
%5. Đọc và phân tích đồ thị quá trình chuyển thể
\begin{hoatdong}{Khởi động}
	\begin{itemize}
		\item Hãy dựa trên những kiến thức đã học về cấu tạo chất để giải thích tại sao cùng một chất lại có thể tồn tại ở các thể khác nhau là rắn, lỏng, khí.
	\end{itemize}
\end{hoatdong}
\subsubsection{Mô hình động học phân tử về cấu tạo chất}

\haicotchay{
	Mô hình động học phân tử về cấu tạo chất có những nội dung cơ bản sau đây:
	\begin{enumerate}
		\item Các chất được cấu tạo từ các hạt riêng biệt là phân tử.
		\item Các phân tử chuyển động không ngừng. Nhiệt độ của vật càng cao thì tốc độ chuyển động của các phân tử cấu tạo nên vật càng lớn.
		\item Giữa các phân tử có lực hút và đẩy gọi chung là lực liên kết phân tử.
	\end{enumerate}
	\chuy{Dùng mô hình này có thể giải thích được cấu trúc của các chất rắn, chất lỏng, chất khí và sự chuyển thể.}
	
}{
	\centering
	\resizebox{\linewidth}{!}{%
		\begin{tikzpicture}
			\draw[thick, fill=blue!5] (0,0) rectangle (4,4);
			\draw[fill=blue] (1,1) circle (0.1);
			\draw[->, thick] (1,1) -- (1.5,1.5);
			\draw[fill=blue] (2,3) circle (0.1);
			\draw[->, thick] (2,3) -- (1.5,2.5);
			\draw[fill=blue] (3,1.5) circle (0.1);
			\draw[->, thick] (3,1.5) -- (3.5,1.2);
			\draw[fill=orange] (2,2) circle (0.3) node[below] {Phấn hoa};
			\draw[->, ultra thick, orange] (2,2) -- (2.3,1.6);
	
		\end{tikzpicture}%
	}
		\captionof{figure}{Chuyển động của phân tử}
	\label{fig: Chuyển động của phân tử  }
}
\begin{emcobiet}
	\begin{itemize}
		\item Thuật ngữ phân tử trong mô hình trên được dùng để chỉ chung các hạt cấu tạo chất: phân tử, nguyên tử, ion.
		\item Người ta thường gọi chuyển động hỗn loạn, không ngừng của các hạt rất nhỏ dưới tác dụng của các phân tử chất lỏng hoặc chất khí là chuyển động Brown.
	\end{itemize}
\end{emcobiet}
\begin{hoatdong}{Tìm hiểu lịch sử và hiện tượng thực tế}
	\begin{itemize}
		\item Trong lịch sử phát triển của khoa học, có hai quan điểm khác nhau về cấu tạo chất là quan điểm chất có cấu tạo liên tục và chất có cấu tạo gián đoạn. Mô hình động học phân tử được xây dựng trên quan điểm nào?
		\item Năm 1827, khi làm thí nghiệm quan sát các hạt phấn hoa rất nhỏ trong nước bằng kính hiển vi, Brown thấy chúng chuyển động hỗn loạn, không ngừng. Chuyển động này được gọi là chuyển động Brown.
		\begin{itemize}
			\item Tại sao thí nghiệm của Brown được coi là một trong những thí nghiệm chứng tỏ các phân tử chuyển động hỗn loạn, không ngừng?
			\item Làm thế nào để với thí nghiệm của Brown có thể chứng tỏ được khi nhiệt độ của nước càng cao thì các phân tử nước chuyển động càng nhanh?
		\end{itemize}
		\item Hãy tìm các hiện tượng thực tế để chứng tỏ giữa các phân tử có lực đẩy, lực hút.
	\end{itemize}
\end{hoatdong}
\subsubsection{Cấu trúc của chất rắn, chất lỏng và chất khí}

 Dựa vào các đặc điểm sau đây của phân tử có thể nêu được sơ lược cấu trúc của hầu hết các chất rắn, chất lỏng, chất khí
\begin{itemize}
	\item Khoảng cách giữa các phân tử càng lớn thì lực liên kết giữa chúng càng yếu.
	\item Các phân tử sắp xếp có trật tự thì lực liên kết giữa chúng mạnh.
\end{itemize}

		\begin{hoatdong}{Mô tả và so sánh các thể của chất}
		\begin{enumerate}
			\item Hãy dựa vào Hình \ref{fig:three-states-v2} để mô tả, so sánh khoảng cách và sự sắp xếp (a), chuyển động (b) của phân tử ở các thể khác nhau. Từ đó mô tả một cách sơ lược về cấu trúc của chất rắn, chất lỏng, chất khí.
			\centering
			\resizebox{\linewidth}{!}{%
				\begin{tikzpicture}[>=stealth, thick, scale=0.9]
					
					%================================================
					% THỂ RẮN (Solid)
					%================================================
					\begin{scope}[shift={(0,0)}]
						% Khung hình
						\draw[line width=2pt] (0,0) rectangle (2.5,2.2);
						\draw[line width=1pt] (2.7,0) rectangle (4.2,2.2);
						
						% Nhãn
						\node[below, font=\small\bfseries] at (1.25,-0.35) {a)};
						\node[below, font=\small\bfseries] at (3.45,-0.35) {b)};
						\node[below, font=\Large\bfseries, color=blue] at (1.75,-0.8) {Thể rắn};
						
						% Phần a) - Cấu trúc rắn
						\foreach \x in {0.2,0.5,0.8,1.1,1.4,1.7,2.0,2.3}{
							\foreach \y in {0.2,0.5,0.8,1.1,1.4,1.7,2.0}{
								\shade[ball color=blue!70] (\x,\y) circle(0.11cm);
							}
						}
						
						% Phần b) - Dao động tại chỗ
						\shade[ball color=blue!70] (3.0,1.1) circle(0.11cm);
						\shade[ball color=blue!70] (3.95,0.6) circle(0.11cm);
						\shade[ball color=blue!70] (3.5,1.8) circle(0.11cm);
						
						% Vẽ mũi tên dao động
						\foreach \x/\y in {3.0/1.1, 3.95/0.6, 3.5/1.8}{
							\draw[<->, dotted, line width=0.8pt] (\x-0.3,\y) -- (\x+0.3,\y);
							\draw[<->, dotted, line width=0.8pt] (\x,\y-0.3) -- (\x,\y+0.3);
						}
					\end{scope}
					
					
					%================================================
					% THỂ LỎNG (Liquid)
					%================================================
					\begin{scope}[shift={(5,0)}]
						% Khung hình
						\draw[line width=2pt] (0,0) rectangle (2.5,2.2);
						\draw[line width=1pt] (2.7,0) rectangle (4.2,2.2);
						
						% Nhãn
						\node[below, font=\small\bfseries] at (1.25,-0.35) {a)};
						\node[below, font=\small\bfseries] at (3.45,-0.35) {b)};
						\node[below, font=\Large\bfseries, color=blue] at (1.75,-0.8) {Thể lỏng};
						
						% Phần a) - Cấu trúc lỏng (gần nhau nhưng không trật tự)
						\foreach \x/\y in {
							0.25/0.25, 0.55/0.45, 0.85/0.2, 1.15/0.55, 1.45/0.3, 1.75/0.5, 2.05/0.25, 2.35/0.4,
							0.35/0.75, 0.7/0.9, 1.0/0.7, 1.3/0.95, 1.6/0.8, 1.9/0.95, 2.2/0.75,
							0.2/1.2, 0.55/1.35, 0.9/1.15, 1.25/1.4, 1.6/1.2, 1.95/1.35, 2.3/1.25,
							0.4/1.65, 0.8/1.8, 1.15/1.65, 1.5/1.85, 1.85/1.7, 2.2/1.85,
							0.3/2.0, 0.75/2.1, 1.2/2.0, 1.65/2.1, 2.1/2.05
						}{
							\shade[ball color=blue!70] (\x,\y) circle(0.11cm);
						}
						
						% Phần b) - Chuyển động (mũi tên hỗn loạn)
						\shade[ball color=blue!70] (2.95,1.8) circle(0.11cm);
						\shade[ball color=blue!70] (3.65,1.2) circle(0.11cm);
						\shade[ball color=blue!70] (3.2,0.5) circle(0.11cm);
						\shade[ball color=blue!70] (3.95,0.8) circle(0.11cm);
						
						\draw[->, dotted, line width=0.8pt] (2.95,1.8) -- (3.65,1.2);
						\draw[->, dotted, line width=0.8pt] (3.65,1.2) -- (3.2,0.5);
						\draw[->, dotted, line width=0.8pt] (3.2,0.5) -- (3.95,0.8);
						\draw[->, dotted, line width=0.8pt] (3.95,0.8) -- (2.95,1.8);
					\end{scope}
					
					
					%================================================
					% THỂ KHÍ (Gas)
					%================================================
					\begin{scope}[shift={(10.0,0)}]
						% Khung hình
						\draw[line width=2pt] (0,0) rectangle (2.5,2.2);
						\draw[line width=1pt] (2.7,0) rectangle (4.2,2.2);
						
						% Nhãn
						\node[below, font=\small\bfseries] at (1.25,-0.35) {a)};
						\node[below, font=\small\bfseries] at (3.45,-0.35) {b)};
						\node[below, font=\Large\bfseries, color=blue] at (1.75,-0.8) {Thể khí};
						
						% Phần a) - Phân tử xa nhau
						\foreach \x/\y in {
							0.4/0.3, 0.9/1.5, 1.6/0.8, 2.1/1.9, 2.3/0.5
						}{
							\shade[ball color=blue!70] (\x,\y) circle(0.11cm);
						}
						
						% Phần b) - Chuyển động nhanh hỗn loạn
						\shade[ball color=blue!70] (2.95,2.0) circle(0.11cm);
						\shade[ball color=blue!70] (3.7,1.0) circle(0.11cm);
						\shade[ball color=blue!70] (3.1,0.4) circle(0.11cm);
						
						\draw[->, dotted, line width=0.8pt] (2.95,2.0) -- (3.7,1.0);
						\draw[->, dotted, line width=0.8pt] (3.7,1.0) -- (3.1,0.4);
						\draw[->, dotted, line width=0.8pt] (3.1,0.4) -- (2.95,2.0);
						\draw[->, dotted, line width=0.8pt] (3.7,1.0) -- (3.95,1.6);
					\end{scope}
					
				\end{tikzpicture}
			}
			\captionof{figure}{Thể rắn -- Thể lỏng -- Thể khí}
			\label{fig:three-states-v2}
			
			\item Hãy giải thích các đặc điểm sau đây của thể khí, thể rắn, thể lỏng:
			\begin{enumerate}
				\item Chất khí không có hình dạng và thể tích riêng, luôn chiếm toàn bộ thể tích bình chứa và có thể nén được dễ dàng.
				\item Vật ở thể rắn có thể tích và hình dạng riêng, rất khó nén.
				\item Vật ở thể lỏng có thể tích riêng nhưng không có hình dạng riêng.
			\end{enumerate}
		\end{enumerate}
	\end{hoatdong}

\subsubsection{Sự chuyển thể}
\paragraph{Khái niệm sự chuyển thể}
\haicotchay{
	\begin{itemize}
		\item Các chất có thể chuyển từ thể này sang thể khác.
		\item Đa số các chất ở thể rắn khi nóng lên có thể chuyển sang thể lỏng, rồi từ thể lỏng sang thể khí. Ngược lại, đa số chất khí khi lạnh đi có thể chuyển sang thể lỏng, rồi từ thể lỏng sang thể rắn.
		\item Một số chất có thể chuyển trực tiếp từ thể rắn sang thể khí và ngược lại. (\ref{fig:Sơ đồ các hình thức chuyển thể})
	\end{itemize}
}{
	\centering
	\resizebox{0.82\textwidth}{!}{
	\begin{tikzpicture}[>=Stealth, thick, font=\sffamily]
		
		\coordinate (RAN) at (0,0);
		\coordinate (LONG) at (5,0);
		\coordinate (KHI) at ($(RAN)!0.5!(LONG)+(0,4.33)$); % tam giác đều
		
		\shade[ball color=green!50!white] (RAN) circle (1.1cm);
		\node at (RAN) {\textbf{RẮN}};
		
		\shade[ball color=red!50!white] (LONG) circle (1.1cm);
		\node at (LONG) {\textbf{LỎNG}};
		
		\shade[ball color=cyan!50!white] (KHI) circle (1.1cm);
		\node at (KHI) {\textbf{KHÍ}};
		
		\draw[->, thick, green!60!black]
		($(RAN)+(1.1,0.3)$) -- ($(LONG)+(-1.1,0.3)$)
		node[midway, above, font=\small\bfseries, text=green!60!black] {nóng chảy};
		
		\draw[->, thick]
		($(LONG)+(-1.25,0)$) -- ($(RAN)+(1.25,0)$)
		node[midway, below, font=\small\bfseries] {đông đặc};
		
		\draw[->, thick, red!70!black]
		($(LONG)+(-0.6,1.0)$) -- ($(KHI)+(0.6,-1.0)$)
		node[midway, above, rotate=-60, font=\small\bfseries, text=red!70!black] {hoá hơi};
		
		\draw[->, thick]
		($(KHI)+(0.6,-1.4)$) -- ($(LONG)+(-0.8,0.9)$)
		node[midway, below, rotate=-65, font=\small\bfseries] {ngưng tụ};
		
		\draw[->, thick]
		($(RAN)+(0.6,1)$) -- ($(KHI)+(-0.7,-1.2)$)
		node[midway, above, rotate=60, font=\small\bfseries] {thăng hoa};
		
		\draw[->, thick, cyan!60!black]
		($(KHI)+(-0.6,-1.4)$) -- ($(RAN)+(0.8,0.9)$)
		node[midway, below, rotate=60, font=\small\bfseries, text=cyan!60!black] {ngưng kết};
		
	\end{tikzpicture}
}
	\captionof{figure}{Sơ đồ các hình thức chuyển thể}
\label{fig:Sơ đồ các hình thức chuyển thể}
	}

\paragraph{Dùng mô hình động học phân tử giải thích sự chuyển thể}
\begin{itemize}
	\item Dựa vào đặc điểm của mô hình động học phân tử có thể giải thích sơ lược các hiện tượng liên quan đến sự chuyển thể như sự hóa hơi, sự nóng chảy.
	\item Trong khi chuyển động hỗn loạn, các phân tử có thể va chạm vào nhau, truyền năng lượng cho nhau. Càng nhận được nhiều năng lượng thì các phân tử chuyển động hỗn loạn càng nhanh, khoảng cách trung bình giữa chúng càng tăng, lực liên kết giữa chúng càng yếu.
\end{itemize}
	\subparagraph{Giải thích sự hóa hơi}
\begin{itemize}
	\item Sự hóa hơi có thể xảy ra dưới hai hình thức là bay hơi và sôi.
	\item \indam{ Bay hơi} là sự hóa hơi xảy ra ở mặt thoáng của chất lỏng. Nước đựng trong một cốc không đậy kín cạn dần là một ví dụ về sự bay hơi.
	\item  \indam{Giải thích cơ chế bay hơi}: Do các phân tử chuyển động hỗn loạn có thể va chạm vào nhau, truyền năng lượng cho nhau nên có một số phân tử ở gần mặt thoáng của chất lỏng có thể có động năng đủ lớn để thắng lực liên kết của các phân tử chất lỏng khác, thoát được ra khỏi mặt thoáng của chất lỏng trở thành các phân tử ở thể hơi.
	\item \indam{Sự sôi}: Khi nhiệt độ của nước tăng tới khoảng $100^\circ\text{C}$ dưới áp suất tiêu chuẩn ($1,01 \cdot 10^5\text{ Pa}$) thì nước bắt đầu sôi. Khi đó các bọt chứa không khí và hơi nước nổi lên trên lòng nước ngày càng nhiều, càng nổi lên trên thể tích các bọt khí này càng tăng, tới mặt thoáng thì vỡ, không khí và hơi nước thoát ra ngoài khí quyển. Do đó sự sôi là sự hóa hơi xảy ra đồng thời ở bên trong và trên mặt thoáng của chất lỏng.
\end{itemize}
\haicotchay{

\begin{hoatdong}{Giải thích hiện tượng hóa hơi}
	\chuy{Tại sao khi bay hơi nhiệt độ của chất lỏng giảm?}
	\begin{enumerate}
		\item Hãy dựa vào đồ thị ở Hình \ref{fig:Đồ thị về sự trao đổi nhiệt độ của nước theo thời gian khi được đun sôi} để mô tả sự thay đổi nhiệt độ của nước khi được đun từ $20^\circ\text{C}$ tới khi sôi.
		\item Khi nước đang sôi thì năng lượng mà nước nhận được từ nguồn nhiệt có được chuyển hóa thành động năng của các phân tử nước không? Tại sao?
	\end{enumerate}
	
	
	\end{hoatdong}}{		\centering
	\resizebox{\linewidth}{!}{%
		\begin{tikzpicture}
			\draw[->] (0,0) -- (4,0) node[right] {Thời gian (s)};
			\draw[->] (0,0) -- (0,3) node[above] {Nhiệt độ ($^\circ\text{C}$)};
			\draw[thick, blue] (0,0.5) node[left]{20} node[right] {(A)} -- (1.5,2.2) node[above] { (B)} -- (3.5,2.2) node[right] {(C)};
			\draw[dashed] (1.5,0) -- (1.5,2.2)  (1.5,2.2)--(0,2.2) node[left, blue]{100};
			\draw[dashed] (3.5,0) -- (3.5,2.2);
			\node at (1.5,-0.3) {$t_1$};
			\node at (3.5,-0.3) {$t_2$};
		\end{tikzpicture}%
	}
	\captionof{figure}{Đồ thị về sự trao đổi nhiệt độ của nước theo thời gian khi được đun sôi}
	\label{fig:Đồ thị về sự trao đổi nhiệt độ của nước theo thời gian khi được đun sôi}
}

\begin{emcobiet}
\haicotchay{
	Chất rắn kết tinh và chất rắn vô định hình:
	\begin{enumerate}
		\item Chất rắn kết tinh là chất rắn cấu tạo nên từ các hạt (phân tử, nguyên tử, ion) cấu tạo nên nó ở thể rắn, liên kết với nhau một cách chặt chẽ, sắp xếp theo một trật tự hình học tuần hoàn tạo thành các mạng tinh thể.
		\item Chất rắn vô định hình là chất ở thể rắn mà các hạt tạo nên nó không tạo thành mạng tinh thể.
		\item Hầu hết các kim loại là các chất kết tinh; thủy tinh, nhựa đường, các chất dẻo,... là các chất vô định hình.
		\item Mỗi chất rắn kết tinh có nhiệt độ nóng chảy (hoặc đông đặc) riêng xác định.
	\end{enumerate}}{

	\centering
	\begin{tikzpicture}[
		x={(0.82cm,0cm)},
		y={(0cm,0.82cm)},
		z={(-0.55cm,0.45cm)},
		line cap=round,
		line join=round,
		>=Stealth,
		font=\sffamily
		]
		
		%================ MÀU =================
		\definecolor{mauIon}{RGB}{0,130,190}
		
		%================ THÔNG SỐ =================
		\def\n{3}
		\def\rNa{2.4pt}
		\def\rCl{4.2pt}
		
		%================ VẼ MẠNG LƯỚI =================
		% Các đường theo phương x
		\foreach \y in {0,1,2,3}{
			\foreach \z in {0,1,2,3}{
				\draw[mauIon!80,line width=0.55pt] (0,\y,\z) -- (3,\y,\z);
			}
		}
		
		% Các đường theo phương y
		\foreach \x in {0,1,2,3}{
			\foreach \z in {0,1,2,3}{
				\draw[mauIon!80,line width=0.55pt] (\x,0,\z) -- (\x,3,\z);
			}
		}
		
		% Các đường theo phương z
		\foreach \x in {0,1,2,3}{
			\foreach \y in {0,1,2,3}{
				\draw[mauIon!80,line width=0.55pt] (\x,\y,0) -- (\x,\y,3);
			}
		}
		
		% Một số đường khuất nét đứt bên trong
		\draw[mauIon!70,dashed,line width=0.55pt] (0,1,1) -- (3,1,1);
		\draw[mauIon!70,dashed,line width=0.55pt] (0,2,2) -- (3,2,2);
		\draw[mauIon!70,dashed,line width=0.55pt] (1,0,1) -- (1,3,1);
		\draw[mauIon!70,dashed,line width=0.55pt] (2,0,2) -- (2,3,2);
		\draw[mauIon!70,dashed,line width=0.55pt] (1,1,0) -- (1,1,3);
		\draw[mauIon!70,dashed,line width=0.55pt] (2,2,0) -- (2,2,3);
		
		%================ KHUNG NGOÀI ĐẬM =================
		\draw[mauIon,line width=1.05pt] (0,0,0) -- (3,0,0) -- (3,3,0) -- (0,3,0) -- cycle;
		\draw[mauIon,line width=1.05pt] (0,0,3) -- (3,0,3) -- (3,3,3) -- (0,3,3) -- cycle;
		\draw[mauIon,line width=1.05pt] (0,0,0) -- (0,0,3);
		\draw[mauIon,line width=1.05pt] (3,0,0) -- (3,0,3);
		\draw[mauIon,line width=1.05pt] (3,3,0) -- (3,3,3);
		\draw[mauIon,line width=1.05pt] (0,3,0) -- (0,3,3);
		
		%================ CÁC ION =================
		% Vẽ từ sau ra trước để hình có chiều sâu
		\foreach \z in {3,2,1,0}{
			\foreach \y in {0,1,2,3}{
				\foreach \x in {0,1,2,3}{
					\pgfmathtruncatemacro{\loai}{mod(\x+\y+\z,2)}
					\ifnum\loai=0
					% Na+ : chấm xanh đặc
					\filldraw[mauIon,draw=mauIon,line width=0.4pt]
					(\x,\y,\z) circle[radius=\rNa];
					\else
					% Cl- : vòng tròn trắng
					\filldraw[fill=white,draw=mauIon,line width=0.85pt]
					(\x,\y,\z) circle[radius=\rCl];
					\fi
				}
			}
		}
		
		%================ NHÃN Na+ VÀ Cl- =================
		\node[font=\scriptsize\bfseries] (Na) at (2.45,3.85,0) {$\mathrm{Na}^{+}$};
		\draw[->,black,line width=0.55pt] (Na.south) -- (3,3,0);
		
		\node[font=\scriptsize\bfseries] (Cl) at (3.75,3.20,0) {$\mathrm{Cl}^{-}$};
		\draw[->,black,line width=0.55pt] (Cl.west) -- (3,3,1);
		
	\end{tikzpicture}
		\captionof{figure}{Tinh thể muối ăn}
	\label{fig:tinh-the-muoi-an}
}
\end{emcobiet}
\subparagraph{Giải thích sự nóng chảy của chất rắn kết tinh}
\begin{itemize}
	\item Khi bị nung nóng, động năng của các hạt trong mạng tinh thể tăng lên, chuyển động nhiệt của các hạt xung quanh vị trí cân bằng của chúng trở nên mạnh mẽ hơn.
	\item Khi đạt tới nhiệt độ nóng chảy, các hạt có năng lượng đủ lớn để thắng lực liên kết và thoát khỏi vị trí cân bằng, trật tự tinh thể bị phá vỡ hoàn toàn, chất rắn kết tinh chuyển thành chất lỏng.
\end{itemize}

	\haicotchay{
		\begin{hoatdong}{Giải thích hiện tượng nóng chảy}
		\begin{enumerate}
			\item Tại sao chất rắn kết tinh khi được đun nóng có thể chuyển thành chất lỏng?
			\item Hãy dựa vào Hình \ref{fig:Đồ thị sự thay đổi nhiệt độ của chất rắn kết tinh khi được làm nóng chảy} để 
				\begin{enumerate}
			\item 	mô tả quá trình nóng chảy của chất rắn kết tinh.
			\item Giải thích tại sao khi đang nóng chảy, nhiệt độ của chất rắn kết tinh không tăng dù vẫn nhận được nhiệt năng. Năng lượng mà chất rắn kết tinh nhận được lúc này dùng để làm gì?
				\end{enumerate}
		\end{enumerate}
	\end{hoatdong}
	}{
			\centering
					\resizebox{\linewidth}{!}{%
			\begin{tikzpicture}[
				>=Stealth,
				line cap=round,
				line join=round,
				scale=0.95,
				font=\small
				]
				
				%----------------------------------------------------
				% Trục tọa độ
				%----------------------------------------------------
				\draw[->,thick] (0,0) -- (7.8,0)
				node[below right]{Thời gian (s)};
				
				\draw[->,thick] (0,0) -- (0,4.6)
				node[above]{Nhiệt độ ($^\circ$C)};
				
				\node[below left] at (0,0) {$O$};
				
				%----------------------------------------------------
				% Nhiệt độ nóng chảy
				%----------------------------------------------------
				\def\ytc{2.25}
				
				\draw[densely dashed]
				(0,\ytc)--(1.18,\ytc);
				
				\node[left] at (0,\ytc) {$t_c$};
				
				\node[align=left,anchor=west]
				at (0.15,2.65)
				{Nhiệt độ nóng chảy};
				
				%----------------------------------------------------
				% Các mốc thời gian
				%----------------------------------------------------
				\coordinate (A0) at (1.18,0);
				\coordinate (A)  at (1.18,0.95);
				
				\coordinate (B)  at (2.00,\ytc);
				
				\coordinate (C)  at (5.05,\ytc);
				
				\coordinate (D)  at (5.95,3.15);
				
				%----------------------------------------------------
				% Đồ thị
				%----------------------------------------------------
				\draw[very thick]
				(A)--(B)--(C)--(D);
				
				%----------------------------------------------------
				% Đường nét đứt
				%----------------------------------------------------
				\draw[densely dashed] (A0)--(A);
				
				\draw[densely dashed]
				(B)--(2,0);
				
				\draw[densely dashed]
				(C)--(5.05,0);
				
				\draw[densely dashed]
				(D)--(5.95,0);
				
				%----------------------------------------------------
				% Ký hiệu a,b,c
				%----------------------------------------------------
				\node at (1.55,1.65) {a};
				
				\node at (3.55,2.00) {b};
				
				\node at (5.45,2.00) {c};
				
			\end{tikzpicture}
		}
				\captionof{figure}{Đồ thị sự thay đổi nhiệt độ của chất rắn kết tinh khi được làm nóng chảy}
			\label{fig:Đồ thị sự thay đổi nhiệt độ của chất rắn kết tinh khi được làm nóng chảy}
			
			\vspace{2mm}
			
			\small
			\begin{tabular}{@{}l@{}}
				Giai đoạn a: Chất rắn chưa nóng chảy;\\
				Giai đoạn b: Chất rắn đang nóng chảy;\\
				Giai đoạn c: Chất rắn đã nóng chảy\\ hoàn toàn.
			\end{tabular}
	}

\begin{emdahoc}
	\begin{itemize}
		\item Mô hình động học phân tử về cấu tạo chất:
		\begin{itemize}
			\item Các chất được cấu tạo từ các hạt riêng biệt là phân tử.
			\item Các phân tử chuyển động không ngừng. Nhiệt độ của vật càng cao thì tốc độ trung bình của các phân tử càng lớn.
			\item Giữa các phân tử có lực liên kết phân tử.
		\end{itemize}
		\item Các chất có thể chuyển từ thể này sang thể khác.
	\end{itemize}
\end{emdahoc}
\begin{emcobiet}
	Trạng thái Plasma:
	\begin{itemize}
		\item Ở nhiệt độ hàng triệu độ, chất không tồn tại ở các thể rắn, lỏng, khí mà ở một thể đặc biệt gọi là Plasma.
		\item Ở thể này, chất không tồn tại dưới dạng các nguyên tử, phân tử mà dưới dạng các ion.
		\item Plasma rất ít gặp trên Trái Đất nhưng trong vũ trụ có tới trên 99\% các chất tồn tại ở thể Plasma.
	\end{itemize}
\end{emcobiet}
\begin{emcothe}
	\begin{itemize}
		\item Dùng mô hình động học phân tử giải thích được một số hiện tượng liên quan đến sự chuyển thể của các chất.
		\item Tìm hiểu và trình bày được vai trò của sự chuyển thể đối với cuộc sống con người như vòng tuần hoàn nước, công nghệ đúc,...
	\end{itemize}
\end{emcothe}
%=====================================================%
